\documentclass{article}

\usepackage{booktabs}
\usepackage{tabularx}
\usepackage{hyperref}

\hypersetup{
    colorlinks=true,       % false: boxed links; true: colored links
    linkcolor=red,          % color of internal links (change box color with linkbordercolor)
    citecolor=green,        % color of links to bibliography
    filecolor=magenta,      % color of file links
    urlcolor=cyan           % color of external links
}

\title{Hazard Analysis\\\progname}

\author{\authname}

\date{}

\input{../Comments.text}
\input{../Common.text}

\begin{document}

\maketitle
\thispagestyle{empty}

~\newpage

\pagenumbering{roman}

\begin{table}[hp]
\caption{Revision History} \label{TblRevisionHistory}
\begin{tabularx}{\textwidth}{llX}
\toprule
\textbf{Date} & \textbf{Developer(s)} & \textbf{Change}\\
\midrule
Date1 & Name(s) & Description of changes\\
Date2 & Name(s) & Description of changes\\
... & ... & ...\\
\bottomrule
\end{tabularx}
\end{table}

~\newpage

\tableofcontents

~\newpage

\pagenumbering{arabic}

\wss{You are free to modify this template.}

\section{Introduction}

\wss{You can include your definition of what a hazard is here.}

\section{Scope and Purpose of Hazard Analysis}

The hazard analysis document outlines the hazards that may occur while using Flexin'.
The high level hazards in scope of Flexin' include:
\begin{itemize}
    \item \textbf{Physical/Psychological Harm}: Includes injuries that may occur during workouts that users intend to track.
    Also includes any impacts on mental health caused by the platform, such as addiction (similar to Instagram or TikTok) or feelings of embarrassment/shame for missing goals.
    \item \textbf{Financial Loss}: It is imperative to note that any interactions with financial markets on Flexin' are not indicative of real market results.
    Users may take strategies learned on the platform and incur serious financial loss.
    \item \textbf{Data Mismanagement}: Poor management of user data from database or cybersecurity perspectives may result in data leaks or loss of data.
    In addition, users improperly providing data can cause over/under rewarding of virtual currency, and can therefore lead to poor learning.
\end{itemize}

The purpose of the hazard analysis is to identify and mitigate risks that could affect the reliability, correctness and safety of the Flexin' app.
Conduction of this analysis will allow the team to anticipate failures early, and design around the hazards in a way that prioritizes the safety of the users.
The findings from this process will inform safety and security requirements, as well as having an impact on compliance and maintainability requirements.
In addition, this work will contribute to the creation of a fun, safe, and intuitive platform for users to learn about compounding and the importance of fitness and financial literacy.

\section{System Boundaries and Components}

\wss{Dividing the system into components will help you brainstorm the hazards.
You shouldn't do a full design of the components, just get a feel for the major
ones.  For projects that involve hardware, the components will typically include
each individual piece of hardware.  If your software will have a database, or an
important library, these are also potential components.}

\section{Critical Assumptions}

\wss{These assumptions that are made about the software or system.  You should
minimize the number of assumptions that remove potential hazards.  For instance,
you could assume a part will never fail, but it is generally better to include
this potential failure mode.}

\section{Failure Mode and Effect Analysis}

\wss{Include your FMEA table here. This is the most important part of this document.}
\wss{The safety requirements in the table do not have to have the prefix SR.
The most important thing is to show traceability to your SRS. You might trace to
requirements you have already written, or you might need to add new
requirements.}
\wss{If no safety requirement can be devised, other mitigation strategies can be
entered in the table, including strategies involving providing additional
documentation, and/or test cases.}

\section{Safety and Security Requirements}

\wss{Newly discovered requirements.  These should also be added to the SRS.  (A
rationale design process how and why to fake it.)}

\section{Roadmap}

\wss{Which safety requirements will be implemented as part of the capstone timeline?
Which requirements will be implemented in the future?}

\newpage{}

\section{Appendix --- Generative AI Use Disclosure}

\wss{ChatGPT (version 5) was used to brainstorm ideas for the template for this disclose section.}

\subsection{AI Tools Used}

\wss{Give the name of the AI tool(s) used and the version.}

\subsection{How and Where Used}

\wss{Explain what the tool(s) were used for.  Examples include: brainstorming,
explaining a technical concept, reviewing the document, generating or modifying
text, generating or modifying diagrams, identifying errors or omissions or
inconsistencies.}

\wss{For each of the uses, identify which parts of the document were affected.}

\wss{You do not need to report every interaction, but you should disclose the 
uses that materially contributed to your work}

\subsection{Verification of AI-Generated Content}

\wss{\begin{itemize}
    \item Describe how the team checked AI-generated or AI-assisted material
    before including it in the document.
    \item In particular, verify factual claims, technical assertions,
    requirements, calculations, references, and other information for which an
    AI system may produce incorrect or fabricated results.
    \item \textbf{The team is responsible for the accuracy, completeness,
consistency, and appropriateness of the submitted document, regardless of
whether material was generated by a human or an AI system.}
\end{itemize}}

\newpage{}

\section{Appendix --- Reflection}

\wss{Not required for CAS 741}

\input{../Reflection.text}

\begin{enumerate}
    \item What went well while writing this deliverable? 
    \item What pain points did you experience during this deliverable, and how
    did you resolve them?
    \item Which of your listed risks had your team thought of before this
    deliverable, and which did you think of while doing this deliverable? For
    the latter ones (ones you thought of while doing the Hazard Analysis), how
    did they come about?
    \item Other than the risk of physical harm (some projects may not have any
    appreciable risks of this form), list at least 2 other types of risk in
    software products. Why are they important to consider?
\end{enumerate}

\end{document}